\begin{ouverture}[college1]{5}{Atomes, molécules et formules}{Partie 1 · Organisation et transformations de la matière}{Que nous apprend une formule chimique ?}
\objectif{Lire les symboles dans un extrait du tableau périodique.}
\objectif{Distinguer atome, molécule et ion.}
\objectif{Compter les atomes dans une formule et dans une collection de molécules.}
\prerequis{Corps pur, mélange et modèle particulaire.}
\end{ouverture}
\pcPage{Le cours}{Atomes, molécules et formules}
\phantomsection\label{ess-5}
\essSection{1}{Des espèces aux entités}
Une espèce chimique est un ensemble d’entités chimiques identiques. Une molécule est un assemblage électriquement neutre d’atomes liés. Le symbole d’un élément comporte une majuscule éventuellement suivie d’une minuscule : C pour carbone, O pour oxygène, H pour hydrogène, N pour azote. Les boules colorées des modèles codent les atomes ; elles ne représentent pas leur couleur réelle.\par
\pcFigure{figures/c05-molecules.pdf}{Les couleurs sont des conventions du dessin : blanc pour H, rouge pour O, gris pour C et bleu pour N.}
\essSection{2}{Lire une formule}
Dans $\ce{H2O}$, l’indice 2 concerne H seulement : une molécule d’eau contient deux atomes d’hydrogène et un atome d’oxygène. L’absence d’indice signifie 1. On rencontre $\ce{O2}$, $\ce{N2}$, $\ce{H2}$, $\ce{CO2}$, $\ce{CH4}$ et $\ce{N2O}$ : dioxygène, diazote, dihydrogène, dioxyde de carbone, méthane et protoxyde d’azote. Les noms ne suffisent pas : vérifier la formule.\par
\clearpage
\essSection{3}{Utiliser un tableau}
Le tableau périodique classe les éléments. Une case donne notamment le nom, le symbole et le numéro atomique $Z$. On peut chercher le symbole d’un nom connu ou retrouver un nom à partir du symbole. On ne confond pas C et Ca, ni Co et $\ce{CO}$ : la typographie change le sens. Le modèle interne de l’atome sera approfondi en 3e.\par
\pcFigure{figures/c05-tableau.pdf}{Extrait de repérage : les cases sélectionnées ne constituent pas une nouvelle classification.}
\essSection{4}{Multiplier et préparer les ions}
$3\,\ce{CO2}$ signifie trois molécules de dioxyde de carbone : trois atomes C et six atomes O au total. Changer un indice change l’espèce ; changer le coefficient change seulement le nombre d’entités. \textbf{Repère de cycle :} un ion est une entité chargée, par exemple $\ce{Na+}$ ou $\ce{Cl-}$. Une solution peut contenir des ions sans contenir des molécules de sel intactes ; leur description détaillée sera reprise en 3e.\par
\begin{pcEncadre}{method}{Un exemple expliqué}Dans $4\,\ce{CH4}$, chaque molécule contient un C et quatre H. On compte donc $4\times1=4$ atomes C et $4\times4=16$ atomes H, soit 20 atomes au total. Écrire $\ce{CH16}$ serait faux : l’indice ne code pas la quantité totale d’un échantillon.\end{pcEncadre}
\pcPage{Expérimenter et retenir}{Atomes, molécules et formules}
\phantomsection\label{activite-5}
\begin{pcEncadre}{method}{Investigation en classe}\textbf{Mission.} Construire des modèles de $\ce{H2O}$, $\ce{CO2}$, $\ce{O2}$ et $\ce{CH4}$ avec des boules ou des jetons légendés. Avant chaque construction, établir un inventaire des atomes. Photographier ou dessiner les modèles, puis échanger les formules avec un autre groupe. \textbf{À expliquer :} pourquoi les modèles représentent la composition, mais pas l’échelle réelle ni tous les mouvements ; pourquoi deux O reliés ne désignent pas deux molécules de dioxygène.\end{pcEncadre}
\pcSousTitre{Ce que je dois savoir expliquer}\begin{questions}
\item Un symbole identifie un élément ; une formule décrit une composition.
\item Un indice s’applique au symbole qui le précède.
\item Un coefficient multiplie toute la formule.
\item Une molécule est neutre ; un ion porte une charge.
\item Le tableau périodique sert à retrouver nom, symbole et numéro atomique.
\end{questions}
\begin{pcEncadre}{attention}{Une erreur à éviter}Ne pas écrire « O est du dioxygène » : O est le symbole de l’élément oxygène ; la molécule de dioxygène a pour formule $\ce{O2}$.\end{pcEncadre}
\pcRepere{cor-activite-5}{Pistes pour l’investigation}\pcGrille
\clearpage\markright{Je démarre}
\rubrique{Je démarre}
\par\noindent\begin{minipage}{\linewidth}
\begin{exercice}[1]{Quatre symboles}
\begin{questions}
\item Donner les symboles de l’hydrogène et de l’oxygène.
\item Quels éléments sont représentés par C et N ?
\end{questions}
\end{exercice}
\end{minipage}
\par\noindent\begin{minipage}{\linewidth}
\begin{exercice}[1]{Lire l’eau}
\begin{questions}
\item Combien d’atomes H et O dans $\ce{H2O}$ ?
\item Combien d’atomes au total ?
\end{questions}
\end{exercice}
\end{minipage}
\par\noindent\begin{minipage}{\linewidth}
\begin{exercice}[1]{Deux molécules différentes}
\begin{questions}
\item Comparer les compositions de $\ce{CO}$ et $\ce{CO2}$.
\item S’agit-il de la même espèce ?
\end{questions}
\end{exercice}
\end{minipage}
\par\noindent\begin{minipage}{\linewidth}
\begin{exercice}[1]{Retrouver un nom}
\begin{questions}
\item Donner les noms de $\ce{N2}$ et $\ce{H2}$.
\item Quelle formule correspond au méthane ?
\end{questions}
\end{exercice}
\end{minipage}
\par\noindent\begin{minipage}{\linewidth}
\begin{exercice}[1]{Repérer une charge}
\begin{questions}
\item Parmi $\ce{Na+}$, $\ce{H2O}$ et $\ce{Cl-}$, lesquels sont des ions ?
\item Quelle charge portent-ils ?
\end{questions}
\end{exercice}
\end{minipage}
\clearpage\markright{Je m’entraîne}
\rubrique{Je m’entraîne}
\par\noindent\begin{minipage}{\linewidth}
\begin{exercice}[2]{Plusieurs molécules}
\begin{questions}
\item Compter C et O dans $5\,\ce{CO2}$.
\item Compter le nombre total d’atomes.
\end{questions}
\end{exercice}
\end{minipage}
\par\noindent\begin{minipage}{\linewidth}
\begin{exercice}[2]{Le protoxyde d’azote}
\begin{questions}
\item Lire la composition d’une molécule de $\ce{N2O}$.
\item Compter les atomes de chaque type dans trois molécules.
\end{questions}
\end{exercice}
\end{minipage}
\par\noindent\begin{minipage}{\linewidth}
\begin{exercice}[2]{Une case du tableau}
\pcFigureEx{figures/ex-case-aluminium.pdf}{Extrait de repérage du tableau périodique.}
\begin{questions}
\item Une case indique « Aluminium, Al, 13 ». Que représentent Al et 13 ?
\item Pourquoi AL n’est-il pas l’écriture correcte ?
\end{questions}
\end{exercice}
\end{minipage}
\par\noindent\begin{minipage}{\linewidth}
\begin{exercice}[2]{Un mélange de modèles}
\pcFigureEx{figures/ex-melange-molecules.pdf}{}
\begin{questions}
\item Une boîte contient quatre molécules $\ce{O2}$ et deux $\ce{N2}$. Combien d’espèces sont présentes ?
\item Compter les atomes O et N.
\end{questions}
\end{exercice}
\end{minipage}
\par\noindent\begin{minipage}{\linewidth}
\begin{exercice}[2]{Corriger une formule}
\pcFigureEx{figures/ex-eau-modele.pdf}{Les lettres identifient les atomes du modèle.}
\begin{questions}
\item Un dessin montre deux atomes H liés à un O. Un élève écrit $\ce{HO2}$. Corriger.
\item Que décrirait la formule écrite ?
\end{questions}
\end{exercice}
\end{minipage}
\clearpage\markright{Je raisonne}
\rubrique{Je raisonne}
\par\noindent\begin{minipage}{\linewidth}
\begin{exercice}[3]{Même inventaire, autre assemblage}
\pcFigureEx{figures/ex-inventaires.pdf}{}
\begin{questions}
\item Comparer les nombres d’atomes dans $2\,\ce{H2O}$ et dans $2\,\ce{H2}+\ce{O2}$.
\item Pourquoi ces écritures ne décrivent-elles pas les mêmes espèces ?
\end{questions}
\end{exercice}
\end{minipage}
\par\noindent\begin{minipage}{\linewidth}
\begin{exercice}[3]{Une collection à construire}
\begin{questions}
\item On possède 6 jetons C et 20 jetons H. Combien de modèles complets de $\ce{CH4}$ peut-on construire au maximum ?
\item Que reste-t-il ?
\end{questions}
\end{exercice}
\end{minipage}
\par\noindent\begin{minipage}{\linewidth}
\begin{exercice}[3]{Une notation ambiguë}
\begin{questions}
\item Pourquoi Co et $\ce{CO}$ n’ont-ils pas le même sens ?
\item Comment éviter la confusion dans un compte rendu ?
\end{questions}
\end{exercice}
\end{minipage}
\par\noindent\begin{minipage}{\linewidth}
\begin{exercice}[3]{Reconstituer une formule}
\begin{questions}
\item Une molécule modèle contient deux C, six H et un O. Proposer sa formule brute.
\item Cette formule permet-elle de connaître un seul assemblage possible ?
\end{questions}
\end{exercice}
\end{minipage}
\par\noindent\begin{minipage}{\linewidth}
\begin{exercice}[3]{Critiquer une représentation}
\begin{questions}
\item Un dessin de gaz présente des boules rouges de taille réelle annoncée « 1 mm ». Identifier le problème.
\item Comment légender correctement le modèle ?
\end{questions}
\end{exercice}
\end{minipage}
\pcPage{Défi du chapitre}{Atomes, molécules et formules}
\phantomsection\label{defi-ex-5}
\pcDefiTitre{5}{Une commande de modèles}
\begin{questions}
\item Il faut construire $8\,\ce{H2O}$ et $5\,\ce{CO2}$. Combien de jetons H, O et C prévoir ?
\item Peut-on annoncer 13 molécules et 39 atomes ?
\item Pourquoi ce nombre d’atomes ne représente-t-il pas une masse en grammes ?
\end{questions}
\pcGrille
